\begin{table}[tbp]
\centering
\caption{Temporal-validation endpoint enrichment by transferred SSL phenotype. The highest-enrichment phenotype is shown for each endpoint, with event counts. Pregnancy-history endpoints use the ever-pregnant stratum as the primary interpretive display to reduce mechanical enrichment from pregnancy exposure; other endpoints use the full analytic cohort. Prevalence-ratio intervals are stratified cluster bootstrap percentile intervals using public-use \texttt{VEST}/\texttt{VECL} design variables with respondent weights retained. Secondary supervised raw-feature and SSL comparisons are reported in Supplementary Table 10.}
\label{tab:enrichment}
\scriptsize
\setlength{\tabcolsep}{2pt}
\begin{tabular*}{\textwidth}{@{\extracolsep{\fill}}>{\raggedright\arraybackslash}p{0.20\textwidth}lcrcc@{}}
\toprule
Endpoint & Analysis & Top phenotype & Events/N & Top \% & PR (95\% CI) \\
\midrule
Adverse proxy & Ever-preg. & P2 & 48/229 & 23.2 & 2.16 (1.60--2.85) \\
Contraceptive vuln. & Full cohort & P1 & 238/2,224 & 10.4 & 1.33 (1.22--1.44) \\
Fertility/loss help & Full cohort & P2 & 44/242 & 18.6 & 1.67 (1.19--2.19) \\
Fecundity limit./infert. & Full cohort & P2 & 94/242 & 41.0 & 1.46 (1.22--1.69) \\
Mistimed/unwtd. & Ever-preg. & P2 & 181/229 & 80.0 & 1.17 (1.09--1.25) \\
\bottomrule
\end{tabular*}
\end{table}
